\subsection{\texorpdfstring{EXAMPLE II: COMPACT ALGEBRAS OF TYPE \(\mathbf{B}_n\) AND \(\mathbf{D}_n\)}{EXAMPLE II: COMPACT ALGEBRAS OF TYPE \textbackslash mathbf\{B\}\_n AND \textbackslash mathbf\{D\}\_n}}

Let V be a finite dimensional real vector space and Q a separating positive quadratic form on V. The orthogonal group associated to Q (Algebra, Chap. IX) is the subgroup \(\mathbf{O}(\mathbf{Q})\) of \(\mathbf{GL}(\mathbf{V})\) consisting of the automorphisms of the real Hilbert space (V, Q); this is a Lie subgroup of \(\mathbf{GL}(\mathbf{V})\) , whose Lie algebra is the subalgebra \(\mathfrak{o}(\mathbf{Q})\) of \(\mathfrak{gl}(\mathbf{V})\) consisting of the endomorphisms x of V such that \(x^{*} = -x\) (Chap. III, §3, no. 10, Cor. 2 of Prop. 37), \(x^{*}\) denoting the adjoint of x relative to Q. Since the group \(\mathbf{O}(\mathbf{Q})\) is compact, \(\mathfrak{o}(\mathbf{Q})\) is thus a compact real Lie algebra. Put \(\mathbf{SO}(\mathbf{Q}) = \mathbf{O}(\mathbf{Q}) \cap \mathbf{SL}(\mathbf{V})\) ; this is a closed subgroup of finite index of \(\mathbf{O}(\mathbf{Q})\) (of index 2 if dim V ≠ 0), hence also with Lie algebra \(\mathfrak{o}(\mathbf{Q})\) .

When \(V = R^{n}\) and Q is the usual quadratic form (for which the canonical basis of \(R^{n}\) is orthonormal), we write \(\mathbf{O}(n, \mathbf{R})\) , \(\mathbf{SO}(n, \mathbf{R})\) , \(\mathfrak{o}(n, \mathbf{R})\) instead of \(\mathbf{O}(\mathbf{Q})\) , \(\mathbf{SO}(\mathbf{Q})\) , \(\mathfrak{o}(\mathbf{Q})\) . The elements of \(\mathbf{O}(n, \mathbf{R})\) (resp. \(\mathfrak{o}(n, \mathbf{R})\) ) are the matrices

\(A \in \mathrm{M}_n(\mathbf{R})\) such that \(A.^t A = I_n\) (resp. \(A = -^t A\) ), which are said to be orthogonal (resp. anti-symmetric).

Let \(\mathrm{V}_{(\mathbf{C})}\) be the complex vector space associated to \(\mathrm{V}\) and let \(\mathrm{Q}_{(\mathbf{C})}\) be the quadratic form on \(\mathrm{V}_{(\mathbf{C})}\) associated to \(\mathbf{Q}\) . Identify \(\mathfrak{gl}(\mathrm{V})_{(\mathbf{C})}\) with \(\mathfrak{gl}(\mathrm{V}_{(\mathbf{C})})\) ; then \(\mathfrak{o}(\mathrm{Q})_{(\mathbf{C})}\) is identified with \(\mathfrak{o}(\mathrm{Q}_{(\mathbf{C})})\) : this is clear since the map \(x \mapsto x^{*} + x\) from \(\mathfrak{gl}(\mathrm{V}_{(\mathbf{C})})\) to itself is \(\mathbf{C}\) -linear. Since \(\mathfrak{o}(\mathrm{Q}_{(\mathbf{C})})\) is of type \(\mathrm{B}_n\) if \(\dim V = 2n + 1\) , \(n \geq 1\) , and of type \(\mathrm{D}_n\) if \(\dim V = 2n\) , \(n \geq 3\) (Chap. VIII, §13, nos. 2 and 4), we deduce:

PROPOSITION 5. Every compact simple real Lie algebra of type \(B_{n}\) , \(n \geq 1\) (resp. of type \(D_{n}\) , \(n \geq 3\) ) is isomorphic to \(\mathfrak{o}(2n + 1, \mathbf{R})\) (resp. \(\mathfrak{o}(2n, \mathbf{R})\) ).

